(1) 由理想气体 $pV=RT$ 和多方约束 $pV^n=K$，得到 $RT=KV^{1-n}$，故 $TV^{n-1}$ 为常量。两端温度比 $T_2/T_1=(V_2/V_1)^{1-n}$，此式与体积变化方向无关。
(2) 积分体积功 $W=\int_{V_1}^{V_2}KV^{-n}dV=K(V_2^{1-n}-V_1^{1-n})/(1-n)=R(T_2-T_1)/(1-n)$。内能是温度的函数，因此 $\Delta U=C_V(T_2-T_1)$，其中 $C_V$ 为一摩尔定容热容。
(3) 第一定律对外做功号约定为 $Q=\Delta U+W$，故 $Q=[C_V+R/(1-n)](T_2-T_1)$，可定义沿此路径的等效热容 $C_{\rm pol}=dQ/dT=C_V+R/(1-n)$。它可能为负，不意味着单相平衡态的 $C_V$ 为负。
(4) $n=0$ 时得到 $C_{\rm pol}=C_V+R=C_p$；$n=\gamma=1+R/C_V$ 时得到 $C_{\rm pol}=0$，符合可逆绝热。$n\to1$ 时温度差也趋零，必须用积分极限得到等温对外功 $W=RT\ln(V_2/V_1)$。
